\documentclass[11pt]{article}
\usepackage{ap-psych}
\title{Unit 1: Biological Bases of Behavior}
\author{Paul Wang}
\date{November 2026}
\docpart{Unit 1 \textperiodcentered\ Study Guide \textperiodcentered\ 15--25\% of MCQ}

\begin{document}
\maketitle

\begin{abstract}
Unit 1 asks how biology and environment combine to produce behavior. It covers heredity and evolution (1.1), the organization of the nervous system (1.2), neurons, neurotransmitters, hormones, and psychoactive drugs (1.3), brain structures and the tools used to study them (1.4), sleep and dreaming (1.5), and sensation across all senses (1.6). Exam items here are mostly application: a symptom or scenario is given and you name the structure, chemical, or process responsible. Each section ends with check questions; the key is on the last page.
\end{abstract}

\scoringline{Check questions: \blank\ / \qtotal \qquad Target: 18+ \qquad Weakest topic: \blank[8em]}

\section{Heredity and environment (1.1)}

Behavior reflects the \textbf{interaction} of nature (genes) and nurture (experience); the exam rarely rewards ``all nature'' or ``all nurture.''

\begin{terms}
Evolutionary psychology & Explains behavior as traits shaped by \textbf{natural selection} because they helped ancestors survive and reproduce (e.g., fear of snakes, preference for sweet foods). \\
Twin studies & Compare \textbf{identical} (monozygotic, \(\sim\)100\% shared genes) with \textbf{fraternal} (dizygotic, \(\sim\)50\%) twins. Greater identical-twin similarity points to genetic influence. \\
Adoption studies & Compare adopted children with biological vs.\ adoptive parents. Resemblance to biological parents suggests genes; to adoptive parents, environment. \\
Heritability & The proportion of \emph{variation among individuals in a group} attributable to genes. It does not describe one person and changes with environment. \\
Epigenetics & Environmental factors switch genes on or off without changing the DNA sequence. \\
Eugenics & A discredited movement that misused heredity claims to justify selective breeding and discrimination. \\
\end{terms}

\begin{trap}
``Heritability of height is .80'' does \emph{not} mean 80\% of your height comes from genes. It means 80\% of the height differences in that population are linked to genetic differences.
\end{trap}

\begin{q}{3}{1.1 Twin studies}{Identical twins reared apart share genes but not environment, isolating genetic influence.}
Which finding most strongly suggests a genetic influence on a trait?
\ch{Fraternal twins raised together are as alike as identical twins raised together.}{Adopted children resemble their adoptive parents.}{Identical twins raised apart are about as similar as identical twins raised together.}{Siblings raised in the same home differ widely.}
\end{q}

\begin{q}{1}{1.1 Heritability}{Heritability describes variation in a group, never a share of one person's trait.}
A trait has a heritability of .60 in a population. This means that
\ch{60\% of the variation in the trait across that population is linked to genetic differences}{60\% of each person's trait comes from genes}{the trait cannot be changed by environment}{60\% of people inherit the trait}
\end{q}

\section{The nervous system (1.2)}

\begin{terms}
Central nervous system (CNS) & Brain and spinal cord. \\
Peripheral nervous system (PNS) & All nerves outside the CNS. Two branches below. \\
Somatic nervous system & Voluntary control of skeletal muscles; carries sensory input. \\
Autonomic nervous system & Controls glands and organs automatically. \\
\quad Sympathetic & \textbf{Arouses}: fight-or-flight. Raises heart rate and blood pressure, dilates pupils, slows digestion, releases glucose. \\
\quad Parasympathetic & \textbf{Calms}: rest-and-digest. Slows heart, constricts pupils, restarts digestion. \\
\end{terms}

\begin{q}{2}{1.2 Autonomic NS}{Racing heart and stopped digestion under threat = sympathetic arousal.}
A hiker sees a bear; her heart races and her digestion slows. Which division is most active?
\chh{Somatic}{Sympathetic}{Parasympathetic}{Central}
\end{q}

\begin{q}{4}{1.2 Somatic NS}{Voluntary skeletal movement is controlled by the somatic nervous system.}
Deciding to raise your hand to answer a question relies most directly on the
\ch{parasympathetic nervous system}{sympathetic nervous system}{autonomic nervous system}{somatic nervous system}
\end{q}

\section{Neurons, chemistry, and drugs (1.3)}

\subsection{Cells and the reflex arc}
\textbf{Neurons} carry messages; \textbf{glial cells} support, nourish, and insulate them (and form myelin). A \textbf{reflex arc} runs \textbf{sensory (afferent) neuron} \(\rightarrow\) \textbf{interneuron} in the spinal cord \(\rightarrow\) \textbf{motor (efferent) neuron}, so you pull your hand from a stove before the brain registers pain.

Structure: \textbf{dendrites} receive \(\rightarrow\) \textbf{cell body (soma)} \(\rightarrow\) \textbf{axon} carries the signal, sped by the \textbf{myelin sheath} \(\rightarrow\) \textbf{terminal buttons} release neurotransmitters into the \textbf{synapse}.

\subsection{Neural firing}
At rest the inside of the neuron is negative (\(\approx -70\) mV): the \textbf{resting potential}. Excitatory input pushes it to \textbf{threshold}; sodium ions rush in (\textbf{depolarization}) and an \textbf{action potential} travels down the axon. Firing is \textbf{all-or-nothing}: a stronger stimulus increases the \emph{rate} of firing, not the size of each impulse. During the \textbf{refractory period} the neuron cannot fire again. Leftover neurotransmitter is reabsorbed by the sending neuron (\textbf{reuptake}).

\begin{compare}{Disorder}{Mechanism}{Result}
Multiple sclerosis & Myelin sheath degenerates. & Slowed, disrupted signals; weakness, vision and coordination problems. \\
Myasthenia gravis & Immune system blocks acetylcholine receptors at muscles. & Muscle weakness and fatigue. \\
\end{compare}

\subsection{Neurotransmitters and hormones}
\begin{terms}
Acetylcholine (ACh) & Muscle movement, learning and memory. Loss linked to Alzheimer's disease. \\
Dopamine & Reward, motivation, movement. Too little: Parkinson's disease. Excess activity: linked to schizophrenia. \\
Serotonin & Mood, sleep, appetite. Low levels linked to depression. \\
Norepinephrine & Alertness and arousal; low levels linked to depressed mood. \\
GABA & Main \textbf{inhibitory} neurotransmitter; low levels linked to anxiety and seizures. \\
Glutamate & Main \textbf{excitatory} neurotransmitter; memory. Excess can overexcite neurons (migraines, seizures). \\
Endorphins & Natural painkillers; ``runner's high.'' \\
Substance P & Transmits pain signals. \\
\end{terms}

\textbf{Hormones} travel through the bloodstream from the \textbf{endocrine system} (slower, longer-lasting than neurotransmitters):
\textbf{adrenaline} (epinephrine; fight-or-flight), \textbf{leptin} (from fat cells; signals fullness), \textbf{ghrelin} (from the stomach; signals hunger), \textbf{melatonin} (pineal gland; sleepiness in darkness), \textbf{oxytocin} (bonding, trust, childbirth, nursing).

\subsection{Psychoactive drugs}
\textbf{Agonists} mimic or increase a neurotransmitter's effect; \textbf{antagonists} block it. \textbf{Reuptake inhibitors} (e.g., SSRIs) leave more neurotransmitter in the synapse.
\begin{terms}
Depressants & Slow neural activity: \textbf{alcohol} (disinhibition, impaired judgment and memory; increases GABA activity). \\
Stimulants & Speed neural activity: caffeine, nicotine, cocaine, methamphetamine (increase dopamine/norepinephrine activity). \\
Hallucinogens & Distort perception: LSD, marijuana (THC). \\
Opioids & Heroin, morphine, fentanyl: act as endorphin agonists; pain relief, euphoria, high overdose risk. \\
Tolerance & Needing more of a drug to get the same effect. \\
Withdrawal & Discomfort and distress after stopping a drug. Often the opposite of the drug's effect. \\
Addiction & Compulsive use despite harm. \\
\end{terms}

\begin{q}{3}{1.3 Neural firing}{All-or-nothing: intensity is coded by firing rate, not impulse size.}
A light touch and a hard press both trigger a sensory neuron. The hard press produces
\ch{larger action potentials}{a longer refractory period}{a higher rate of firing}{a slower action potential}
\end{q}

\begin{q}{1}{1.3 Neurotransmitters}{Muscle weakness + ACh receptor blocking = myasthenia gravis.}
A patient's immune system blocks acetylcholine receptors at the muscles. The patient most likely has
\ch{myasthenia gravis}{multiple sclerosis}{Parkinson's disease}{schizophrenia}
\end{q}

\begin{q}{4}{1.3 Drugs}{A drug that blocks receptors is an antagonist; reduced effect.}
A drug occupies dopamine receptor sites without activating them. It is acting as
\ch{a dopamine agonist}{a reuptake inhibitor}{a stimulant}{a dopamine antagonist}
\end{q}

\begin{q}{2}{1.3 Hormones}{Ghrelin from the stomach signals hunger.}
Which hormone is released by an empty stomach and increases hunger?
\chh{Leptin}{Ghrelin}{Melatonin}{Oxytocin}
\end{q}

\begin{q}{3}{1.3 Drugs}{Needing more drug for the same effect = tolerance.}
After months of daily coffee, Jorge needs three cups to feel as alert as one used to make him. This illustrates
\chh{withdrawal}{addiction}{tolerance}{reuptake}
\end{q}

\section{The brain (1.4)}

\begin{terms}
Brainstem / medulla & Heartbeat, breathing, blood pressure. Damage is often fatal. \\
Reticular activating system & Network through the brainstem controlling arousal and alertness. \\
Cerebellum & Balance, coordination, timing; procedural (skill) memory. \\
Thalamus & Sensory relay station for all senses except smell. \\
Hypothalamus & Hunger, thirst, body temperature, sexual behavior; controls the pituitary. \\
Pituitary gland & ``Master gland''; directs other endocrine glands. \\
Hippocampus & Forms new explicit (conscious) memories. \\
Amygdala & Fear and aggression; emotional memory. \\
Reward center & Dopamine pathways (e.g., nucleus accumbens) activated by pleasure and drugs. \\
Corpus callosum & Band of fibers connecting the two hemispheres. \\
\end{terms}

\subsection{Cerebral cortex}
\begin{compare}{Lobe}{Key regions}{Functions}
Frontal & Prefrontal cortex; motor cortex (rear strip); Broca's area (left) & Planning, judgment, impulse control, personality; voluntary movement; speech production. \\
Parietal & Somatosensory cortex (front strip) & Touch, temperature, body position; spatial processing. \\
Occipital & Visual cortex & Vision. \\
Temporal & Auditory cortex; Wernicke's area (left) & Hearing; language comprehension. \\
\end{compare}

\textbf{Broca's aphasia}: slow, effortful speech; comprehension largely intact. \textbf{Wernicke's aphasia}: fluent but meaningless speech; poor comprehension.

\subsection{Hemispheres and plasticity}
Each hemisphere controls the \emph{opposite} side of the body (contralateral). Language is usually lateralized to the left. In a \textbf{split-brain} patient (corpus callosum cut), an image shown to the \textbf{left visual field} goes to the right hemisphere: the patient cannot \emph{say} what it is but can pick it out with the \emph{left} hand. \textbf{Neuroplasticity} is the brain's ability to reorganize after damage or experience, greatest in childhood.

\subsection{Tools}
\textbf{Case studies} of brain damage (Phineas Gage: frontal-lobe injury changed personality), \textbf{lesioning} (removing or destroying tissue in animals), \textbf{EEG} (electrical activity; excellent timing), \textbf{fMRI} (blood flow shows which areas are active), and other scans.

\begin{q}{2}{1.4 Brain}{Damage to the hippocampus prevents new explicit memories.}
After surgery, a patient remembers his childhood but cannot form new memories of people he meets. The damaged structure is most likely the
\chh{cerebellum}{hippocampus}{medulla}{occipital lobe}
\end{q}

\begin{q}{4}{1.4 Aphasia}{Fluent but meaningless speech with poor comprehension = Wernicke's area (left temporal).}
A stroke patient speaks fluently but her sentences make little sense, and she struggles to understand others. The damage is most likely in
\ch{Broca's area in the left frontal lobe}{the right parietal lobe}{the cerebellum}{Wernicke's area in the left temporal lobe}
\end{q}

\begin{q}{1}{1.4 Split brain}{Left visual field \(\rightarrow\) right hemisphere: can't name it, but left hand (right-hemisphere control) can select it.}
A split-brain patient sees a key flashed to her \emph{left} visual field. She will most likely
\ch{say she saw nothing but pick out the key with her left hand}{say ``key'' and pick it out with her right hand}{say ``key'' but be unable to pick it out}{be unable to identify it with either hand}
\end{q}

\begin{q}{3}{1.4 Tools}{fMRI shows blood flow, i.e., which areas are active during a task.}
To find which brain areas become more active when people view faces, a researcher would best use
\chh{an EEG}{a lesion}{an fMRI}{a case study}
\end{q}

\section{Sleep (1.5)}

The \textbf{circadian rhythm} is the roughly 24-hour biological clock, set by light through the \textbf{suprachiasmatic nucleus} of the hypothalamus, which signals the pineal gland to release melatonin. \textbf{Jet lag} and \textbf{shift work} disrupt it.

\begin{compare}{Stage}{What happens}{Notes}
NREM-1 & Light sleep; drifting off. & Hypnagogic sensations (falling, flashes). \\
NREM-2 & Deeper; \textbf{sleep spindles} (bursts of activity). & Most of total sleep time. \\
NREM-3 & Deep, slow \textbf{delta waves}; hard to wake. & Sleepwalking and night terrors occur here; growth hormone released. Shrinks across the night. \\
REM & Rapid eye movements, vivid dreams, brain active but muscles paralyzed (``paradoxical sleep''). & Periods lengthen across the night. \textbf{REM rebound}: more REM after deprivation. \\
\end{compare}
A full cycle lasts about 90 minutes and repeats four to six times a night.

\textbf{Dream theories}: \textbf{activation-synthesis} (the brainstem fires randomly during REM and the cortex weaves the signals into a story) and \textbf{memory consolidation} (sleep strengthens and stores the day's memories).

\begin{terms}
Insomnia & Persistent trouble falling or staying asleep. \\
Narcolepsy & Sudden, uncontrollable sleep attacks, sometimes collapsing directly into REM. \\
Sleep apnea & Breathing repeatedly stops during sleep; loud snoring, daytime fatigue. \\
REM sleep behavior disorder & Muscle paralysis fails during REM, so people act out dreams. \\
Somnambulism & Sleepwalking during NREM-3. \\
\end{terms}

\begin{q}{2}{1.5 Sleep stages}{Sleepwalking occurs during deep NREM-3, early in the night.}
A child often sleepwalks during the first two hours of the night. This most likely happens during
\chh{REM sleep}{NREM-3 sleep}{NREM-1 sleep}{the hypnagogic state}
\end{q}

\begin{q}{4}{1.5 Sleep disorders}{Acting out dreams = loss of REM muscle paralysis.}
A man punches the air and shouts while dreaming that he is boxing. He most likely has
\chh{narcolepsy}{sleep apnea}{insomnia}{REM sleep behavior disorder}
\end{q}

\begin{q}{1}{1.5 Dream theory}{Activation-synthesis: random brainstem activity interpreted by the cortex.}
``Dreams are the brain's attempt to make sense of random neural activity from the brainstem.'' This describes
\ch{activation-synthesis theory}{memory consolidation theory}{REM rebound}{the circadian rhythm}
\end{q}

\section{Sensation (1.6)}

\subsection{Thresholds}
\textbf{Transduction}: receptors convert stimulus energy into neural signals. \textbf{Absolute threshold}: the minimum stimulation detected 50\% of the time. \textbf{Difference threshold} (just-noticeable difference): the smallest detectable change. \textbf{Weber's law}: the JND is a constant \emph{proportion} of the original stimulus (a 2\% change in weight, not 2 grams).
\textbf{Signal detection theory}: detection depends on the stimulus \emph{and} the person's expectations, motivation, and fatigue (hits, misses, false alarms, correct rejections). \textbf{Sensory adaptation}: reduced sensitivity to a constant stimulus. \textbf{Sensory interaction}: senses influence each other (smell shapes taste). \textbf{Synesthesia}: one sense triggers another (hearing colors).

\subsection{Vision}
Light passes the cornea, pupil, and \textbf{lens} (which changes shape to focus: accommodation) to the \textbf{retina}. \textbf{Rods}: dim light, peripheral vision, no color. \textbf{Cones}: color and fine detail, packed in the \textbf{fovea}. Signals go to bipolar then ganglion cells, whose axons form the \textbf{optic nerve}; where it leaves the eye there are no receptors (\textbf{blind spot}). Nearsightedness: distant objects blurry; farsightedness: near objects blurry.
\begin{compare}{Theory}{Claim}{Explains}
Trichromatic & Three cone types (red, green, blue). & Color mixing; most color blindness (missing a cone type: dichromatism). \\
Opponent-process & Opposing pairs: red--green, blue--yellow, black--white. & \textbf{Afterimages} (stare at green, see red). \\
\end{compare}
\textbf{Prosopagnosia}: inability to recognize faces. \textbf{Blindsight}: responding to visual stimuli without conscious awareness of seeing them.

\subsection{Hearing}
\textbf{Amplitude} \(\rightarrow\) loudness; \textbf{frequency} \(\rightarrow\) pitch. Path: eardrum \(\rightarrow\) ossicles (hammer, anvil, stirrup) \(\rightarrow\) oval window \(\rightarrow\) \textbf{cochlea}, where hair cells on the basilar membrane transduce vibration \(\rightarrow\) auditory nerve. \textbf{Place theory}: pitch depends on \emph{where} the membrane vibrates (high pitches). \textbf{Frequency theory}: pitch depends on the \emph{rate} of neural firing (low pitches); the \textbf{volley principle} covers middle pitches. \textbf{Conduction deafness}: damage to the eardrum or ossicles (mechanical); \textbf{sensorineural deafness}: damage to hair cells or the auditory nerve. Sounds are \textbf{localized} by differences in timing and loudness between the ears.

\subsection{Chemical, body, and skin senses}
\textbf{Taste} (gustation): sweet, salty, sour, bitter, \textbf{umami}; supertasters have more taste buds. \textbf{Smell} (olfaction) bypasses the thalamus and links directly to the amygdala and hippocampus, which is why smells trigger emotional memories. \textbf{Touch}: pressure, warmth, cold, pain. \textbf{Gate-control theory}: a spinal ``gate'' lets pain through or blocks it; rubbing a bumped elbow or distraction can close it. \textbf{Phantom limb}: sensations in a missing limb. \textbf{Vestibular sense}: balance (semicircular canals). \textbf{Kinesthesis}: position and movement of body parts.

\begin{q}{3}{1.6 Weber's law}{JND is a constant proportion: 2\% of 500 g = 10 g.}
If the JND for weight is 2\%, what is the smallest change a person would notice in a 500-gram package?
\chh{2 grams}{5 grams}{10 grams}{20 grams}
\end{q}

\begin{q}{1}{1.6 Vision}{Afterimages in complementary colors are explained by opponent-process theory.}
After staring at a red square, Maya sees a green square on a white wall. This is best explained by
\ch{opponent-process theory}{trichromatic theory}{place theory}{signal detection theory}
\end{q}

\begin{q}{2}{1.6 Signal detection}{Expectations and motivation shift detection: more ``yes'' responses, including false alarms.}
A nurse expecting a patient's call ``hears'' the call bell several times when it did not ring. Signal detection theory calls these
\chh{misses}{false alarms}{hits}{correct rejections}
\end{q}

\begin{q}{4}{1.6 Hearing}{Damage to hair cells or the auditory nerve = sensorineural deafness.}
Years of loud concerts damaged the hair cells in Leo's cochlea. He has
\ch{conduction deafness}{prosopagnosia}{a damaged vestibular sense}{sensorineural deafness}
\end{q}

\begin{q}{3}{1.6 Pain}{Rubbing sends competing large-fiber signals that close the spinal gate.}
Rubbing a stubbed toe reduces its pain. This is best explained by
\chh{sensory adaptation}{Weber's law}{gate-control theory}{kinesthesis}
\end{q}

\makekey
\end{document}
